\documentclass[11pt,a4paper]{article}

\usepackage[T1]{fontenc}
\usepackage[utf8]{inputenc}
\usepackage[english]{babel}
\usepackage{lmodern}
\usepackage{geometry}
\usepackage{amsmath}
\usepackage{booktabs}
\usepackage{enumitem}

\geometry{margin=2.5cm}

\title{Towards a Modular Inference Engine for Language Models}
\author{}
\date{}

\begin{document}

\maketitle

\section{Design}

\subsection{Problem Statement}

Current inference engines, such as \textit{llama.cpp}, provide robust performance but largely rely on tightly coupled software organizations. This approach favors efficiency within a well-defined usage framework, at the cost of limited flexibility when introducing new computation schemes, new hardware targets, or specific distribution strategies.

In this context, our research problem can be formulated as follows: \textit{how can we design an inference engine that preserves competitive performance while enabling explicit and extensible composition of execution pipeline components?}

\subsection{Limitations of Existing Approaches}

The analysis of existing solutions highlights four major limitations.

\begin{enumerate}[leftmargin=*]
	\item \textbf{Structural coupling to the inference pipeline.} Loading, optimization, and execution stages are often intertwined, making it difficult to replace an isolated component.
	\item \textbf{Imposed architecture.} Existing tools constrain the user to adopt a predefined execution scheme, limiting the exploration of algorithmic or system-level variants.
	\item \textbf{High extension cost.} Adding new operators, partitioning policies, or backends frequently requires cross-cutting code modifications.
	\item \textbf{Limited modular testability.} Evaluating an isolated component is difficult when the architecture requires execution of the full pipeline.
\end{enumerate}

\subsection{Design Hypothesis}

We hypothesize that an explicit intermediate representation of computation, in the form of a directed acyclic graph (DAG), constitutes the appropriate abstraction level for decoupling engine responsibilities.

This hypothesis leads to three structuring design choices:

\begin{itemize}[leftmargin=*]
	\item \textbf{Computation abstraction.} The model is projected into a DAG representation independent of the source format.
	\item \textbf{Separation of ``what'' and ``where''.} The graph encodes computation dependencies, while backend modules define execution on a given hardware target.
	\item \textbf{Composable transformations.} Partitioning and optimization policies are described as parameterizable graph transformations.
\end{itemize}

\subsection{Proposed Modular Architecture}

The architecture is organized into four functional layers:

\begin{enumerate}[leftmargin=*]
	\item \textbf{Model ingestion}: loading from a source format (e.g., GGUF).
	\item \textbf{Common IR}: conversion into a unified DAG (operational nodes, data edges).
	\item \textbf{Optimization/partitioning}: application of decomposition and scheduling strategies.
	\item \textbf{Backend execution}: projection of the execution plan onto one or several targets (CPU, CUDA, Metal, etc.).
\end{enumerate}

This decoupling enables the evolution of one layer without impacting the entire system and facilitates controlled comparison of implementation variants.

\subsection{Design Objectives}

The objectives associated with this proposal are the following:

\begin{itemize}[leftmargin=*]
	\item provide an extensible and stable inference interface;
	\item enable explicit partitioning of the computation graph;
	\item reduce the integration cost of new operators and backends;
	\item improve testability and experimental reproducibility.
\end{itemize}

\subsection{Summary}

In summary, the design contribution is not limited to local execution optimization, but instead proposes a software organization oriented toward research and experimentation. The DAG becomes the central artifact from which partitioning and execution decisions are derived, opening the way to a more rigorous comparative evaluation of modular inference strategies.

\section{Implementation}

\subsection{Overview}

The implementation is developed as a set of loosely coupled C/C++ modules integrated with GGML. The choice of modular granularity serves two practical goals: (i) enabling unit-level validation of each component and (ii) facilitating engine evolution without requiring global rewrites.

The main modules are:
\begin{itemize}[leftmargin=*]
	\item \texttt{parser.c/.h}: loading of a GGUF model and metadata extraction;
	\item \texttt{distrib.c/.h}: representation of the partitioning graph and balancing heuristics;
	\item \texttt{dummy\_distrib.c}: partitioning baseline;
	\item \texttt{KaHIP.c}: advanced partitioning interface using KaHIP;
	\item \texttt{parser\_main.c}: experimental entry point.
\end{itemize}

\subsection{Model Parsing Component}

The \texttt{parser} component implements GGUF loading and instantiates a model structure containing the GGML context, the GGUF context, and the hyperparameters relevant to inference.

Functions of interest:
\begin{itemize}[leftmargin=*]
	\item \texttt{model\_load(...)}: file validation and model initialization;
	\item \texttt{run\_single\_token\_inference(...)}: minimal execution for functional verification;
	\item \texttt{run\_batch\_inference(...)}: multi-token execution;
	\item \texttt{save\_logits(...)}: output export for comparison.
\end{itemize}

The GGUF format is used as the canonical input format. In practice, the parser assumes the availability of a coherent header (version, number of tensors, metadata) and correctly named and typed tensor descriptors.

\subsection{GGUF Parser Details}

The parser is designed to (i) validate the file structure, (ii) construct the corresponding GGML tensors, and (iii) produce an operational representation (DAG) exploitable by the partitioning tool.

Main stages:
\begin{enumerate}[leftmargin=*]
	\item Read the GGUF header and verify the magic/version.
	\item For each tensor descriptor, allocate a \texttt{ggml\_tensor} and record its name, shape, and type.
	\item Construct the \textit{name -> tensor} lookup table to retrieve a tensor from alternative names (e.g., \texttt{"weight"} / \texttt{"W"}).
	\item Construct the forward DAG by traversing operator dependencies in the GGML context (alias and usage extraction functions).
\end{enumerate}

Notable functions and implementation details:
\begin{itemize}[leftmargin=*]
	\item \texttt{find\_tensor\_any(...)}: returns a pointer to an existing tensor among several candidate names; useful for robustness against dump variants.
	\item \texttt{load\_required\_tensor\_any(...)}: abstraction method used to verify and load required tensors.
\end{itemize}

Example mapping heuristic: to estimate the memory size of a tensor, we compute
\[
\mathrm{bytes}(T) = \prod_{d\in \mathrm{dims}} d \times \mathrm{sizeof}(\text{type})
\]
and store this value as an attribute later used for communication cost estimation.

\subsection{Construction of the Partitioning Graph}

The partitioning graph is constructed from the execution DAG: each operation becomes a vertex associated with an estimated computation cost, and an edge is created between two vertices if a tensor produced by one is consumed by the other.

Vertex cost computation (computation cost): a simple estimation used in the implementation is
\[
C_{\mathrm{comp}}(v) = \frac{\mathrm{FLOPs}(v)}{\mathrm{peak\_flops}_{\mathrm{target}}},
\]
where $\mathrm{FLOPs}(v)$ is an estimate of the number of floating-point operations required to execute $v$ (for example, $2\times n \times m$ for a matrix-vector multiplication). The parameter $\mathrm{peak\_flops}_{\mathrm{target}}$ is provided by the user or obtained through profiling.

Edge weight computation (communication cost): for an edge carrying a tensor $T$ of size $B$ bytes,
\[
W_{\mathrm{comm}}(e) = \alpha + \frac{B}{\beta},
\]
where $\alpha$ is the fixed latency (overhead) and $\beta$ is the effective bandwidth. These constants may be estimated from micro-benchmarks.

Construction pseudocode:
\begin{verbatim}
for each op v in DAG:
	comp_cost[v] = estimate_flops(v) / peak_flops
	for each output tensor t of v:
		for each consumer u of t:
			add edge (v -> u) with weight = alpha + bytes(t)/bandwidth
\end{verbatim}

\subsection{Partitioning Algorithms}

\paragraph{Greedy baseline (\texttt{dummy\_distrib})}

The implemented baseline is a greedy bin-packing heuristic: vertices are ordered by decreasing cost and assigned to the partition minimizing the increase in the objective $\mathcal{L}$. This strategy is simple but robust for small graphs.

\paragraph{Partitioning with KaHIP}

For KaHIP, the graph is converted into a format accepted by KaHIP (hypergraph/graph depending on the API), and the call is delegated either through the C API or by invoking the executable. The result is a \texttt{vertex\_to\_partition} table that is re-imported and used to annotate the execution plan.

Typical exchange file (graph): edge list with weights and number of vertices, followed by vertex weights.

\subsection{Interface with KaHIP (Details)}

The integration includes two variants: (i) direct linking with the KaHIP library at compile time, and (ii) graph export to disk followed by invocation of the KaHIP binary.

In the ``API'' path, the expected C structures are built and \texttt{kaffpa\_interface(...)} is called (or the appropriate API depending on the version). In the binary path, a text file \texttt{graph.kt} is produced and the following command is executed:
\begin{verbatim}
kahip --input=graph.kt --k=4 --preconfiguration=fast
\end{verbatim}

Then, the output file associating each vertex with a partition is parsed and copied into \texttt{vertex\_to\_partition}.

\subsection{Tests, Benchmarks, and Validation Protocol}

To guarantee reproducibility, we propose a set of unit tests and micro-benchmarks:

\begin{itemize}[leftmargin=*]
	\item \textbf{Parser unit tests}: verify that the number of tensors read matches the metadata and that the main names exist (assertions on \texttt{n\_tensors}, \texttt{hparams}).
	\item \textbf{Logit regression tests}: execute \texttt{run\_single\_token\_inference} on a standard context and compare logits against reference values (relative tolerance, e.g., $10^{-5}$).
	\item \textbf{Benchmarks}: \texttt{test\_perf} provides simulated execution time per vertex (in ms), estimated inter-partition cut, and balance ratio.
\end{itemize}

Minimal commands:
\begin{verbatim}
mkdir -p build && cd build
cmake .. -DGGML_USE_KAHIP=ON -DKAHIP_ROOT=/path/to/kahip-install
cmake --build . --target parser test_perf model_compare
// Run parser
./build/parser sample_model.gguf
// Run perf
./build/test_perf sample_model.gguf
\end{verbatim}

Acceptance criteria:
\begin{itemize}[leftmargin=*]
	\item correct load structure (acyclic DAG);
	\item logit difference between reference and implementation $<$ tolerance;
	\item valid partitions (no unassigned vertex), reasonable balance ratio (e.g., $<1.2$).
\end{itemize}

\subsection{Implementation Notes and Limitations}

Several implementation choices deserve emphasis:
\begin{itemize}[leftmargin=*]
	\item FLOP estimation is approximate and depends on the operation type; runtime instrumentation improves precision;
	\item the interface with KaHIP assumes correct graph conversion into an accepted structure; this may imply information loss regarding edge costs if a simplified representation is used;
	\item the current implementation does not handle dynamic operators or explicit loops; such cases require prior graph transformation.
\end{itemize}

\subsection{Partitioning Component}

The \texttt{distrib} component transforms the computation graph into a weighted partitioning graph: vertices represent the computational cost of operations, while edges represent the communication cost between dependent operations.

The optimized objective is defined as:
\[
\mathcal{L} = \lambda_{\mathrm{comm}}\,C_{\mathrm{cut}} + \lambda_{\mathrm{bal}}\,P_{\mathrm{balance}},
\]
where $C_{\mathrm{cut}}$ is the inter-partition cut cost and $P_{\mathrm{balance}}$ is the load imbalance penalty. This formulation enables explicit adjustment of the communication/balancing trade-off.

The partitioning configuration notably includes: the number of partitions, objective weights, a hard-cap load limit, and iteration parameters for local refinement.

\subsection{KaHIP Integration}

The \texttt{KaHIP.c} module encapsulates the use of KaHIP as an external partitioning solver. Integration is optional at CMake configuration time, allowing comparison between two regimes:
\begin{itemize}[leftmargin=*]
	\item \textbf{internal baseline} (without KaHIP);
	\item \textbf{advanced partitioning} (with KaHIP).
\end{itemize}

This duality provides a useful experimental basis for quantifying the contribution of an external partitioner on graphs of different sizes.

\subsection{Reproducibility Details}

\subsubsection{Validation Protocol}

To reproduce the baseline technical results, we recommend the following sequence:
\begin{enumerate}[leftmargin=*]
	\item verify GGUF model loading via \texttt{parser};
	\item validate output consistency (logits) on a simple case;
	\item execute benchmarks (\texttt{test\_perf});
	\item compare partitioning strategies (\texttt{model\_compare}).
\end{enumerate}

\subsection{Implementation Discussion}

The current implementation constitutes a reproducible foundation for studying inference graph partitioning, but does not yet cover full distributed multi-backend execution. Consequently, the results obtained at this stage should be interpreted as feasibility and characterization results rather than as a complete end-to-end system validation.

\end{document}